% =============================================================================
%  lecture.tex -- lecture notes / skrypt template
%  Compile with: latexmk lecture.tex   (uses lualatex via .latexmkrc)
% =============================================================================
\documentclass[11pt, a4paper]{article}
\usepackage{mypreamble}

% --- Document metadata ---
\newcommand{\coursename}{GAL26}
\newcommand{\notetaker}{Radosław Galiński}
\newcommand{\semester}{Semestr letni 2025/26}

\pagestyle{fancy}
\fancyhf{}
\fancyhead[L]{\itshape\coursename}
\fancyhead[R]{\itshape\nouppercase\leftmark}
\fancyfoot[C]{\thepage}
\renewcommand{\headrulewidth}{0.4pt}

\begin{document}

\thispagestyle{empty}
\lectureheader{\coursename}{\notetaker}{\semester}

\tableofcontents
\bigskip
\newpage

% =============================================================================
\section{Wprowadzenie}
% use gq in vim to style

To są moje notatki z GAL'u z semetru letniego 25/26 na drugie kolokwium. Jeśli
nie ma czegoś lub jest słabo opisane to proszę nie krępować się poprawić.

\section{Iloczyn tensorowy}

\begin{definition}
    Na iloczynie kartezjańskim \(V_1 \times ... \times V_n\) przestrzeni nad
    \(K\) wprowadzamy strukturę przestrzeni liniowej definiując operacje:
    \begin{enumerate}
        \item \((v_1, ..., v_n) + (v_1', ..., v_n') = (v_1 + v_1', ..., v_n + v_n')\);
        \item \(\lambda \cdot (v_1, ..., v_n) = (\lambda \cdot v_1, ..., \lambda \cdot v_n)\).
    \end{enumerate}
\end{definition}
\begin{definition}
    Przekształcenie \(n\)-liniowe \(\phi\) definiujemy jako przekształcenie na iloczynie
    kartezjańskim przestrzeni liniowych \(V_1 \times ... \times V_n\), tak że zachodzi:
    \begin{enumerate}
    \item \(\lambda \cdot \phi(v_1, ..., v_n) = \phi(v_1, ..., \lambda \cdot v_i, ..., v_n)\);
    \item \(\phi(v_1, ..., v_n) + \phi(v_1, ..., v_{i-1}, v_i', v_{i + 1}, ...,
    v_n) = \phi(v_1, ..., v_{i-1}, v_i + v_i', v_{i + 1}, ..., v_n)\).
    \end{enumerate}
\end{definition}
\begin{definition}
    Niech \(V_1, ..., V_n\) będą przestrzeniami liniowymi nad \(K\). Wówczas definiujemy
    iloczyn tensorowy \(T, \tau\) jako parę:
    \begin{enumerate}
        \item przestrzeni liniowej \(T\) oznaczanej przez \(V_1 \otimes ... \otimes V_n \),
        \item przekształcenia \(n\)-liniowego \(\tau: V_1 \times ... \times V_n \to T\),
    \end{enumerate}
    taką, że dla dowolnego przekształcenia \(n\)-liniowego \(\phi\) istnieje
    dokładnie jedno przekształcenie \(f\), że \(\phi = f \circ \tau\).
\end{definition}

\begin{theorem}
	Iloczyn tensorowy \(V_1 \otimes... \otimes V_n\) jest wyznaczony z
	dokładnością do izomorfizmu.
\end{theorem}

\begin{remark}
	Jeśli \((\eps_1, ..., \eps_n) \subseteq U\) to układ liniowo niezależny, to jeśli dla
	pewnego układu wektorów \((v_1, ..., v_n) \subseteq V\) mamy
	\[
	    \eps_1 \otimes v_1 + ... + \eps_n \otimes v_n = 0
	\]
	to \(v_1, ..., v_n = 0\)
\end{remark}
\begin{remark}
	Niech \((\eps_1, ..., \eps_n) \subseteq U\) to układ liniowo niezależny.
	Jeśli \((v_1, ..., v_n) \subseteq V\) spełnia warunek
	\[
	\eps_1 \otimes v_1 + ... + \eps_n \otimes v_n = z_1 \otimes w_1 + ... + z_n \otimes w_n
	\]
	dla pewnych \(z_i, w_i\), to dla każdego \(j\) zachodzi
	\[
	    v_i \in lin(w_1, ..., w_n)
	\]
\end{remark}

\begin{theorem}
	Niech \((\alpha_1, ..., \alpha_n) \subseteq V\) będzie bazą \(V\) oraz
	\((\beta_1, ..., \beta_m)\) będzie bazą \(W\). Wtedy układ
	\[
	\{ \alpha_i \otimes \beta_j \}
	\]
	jest bazą przestrzeni \(V \otimes W\).
\end{theorem}

\begin{definition}
	Definiujemy rząd tensora \(t\) jako minimalną liczbę, że istnieje układ tensorów
	prostych którego kombinacja liniowa jest równa \(t\).
\end{definition}
\section{Iloczyn zewnętrzny}
\begin{definition}
    Przekształcenie \(n\)-liniowe \(f\) jest:
    \begin{enumerate}
        \item symetryczne, jeśli dla dowolnej permutacji \(\sigma\) jest
        \[f(v_1, ..., v_n) = f(v_{\sigma(1)}, ..., v_{\sigma(n)});\]

        \item antysymetryczne, jeśli dla dowolnej permutacji \(\sigma\) jest
        \[(f(v_1, ..., v_n) = \text{sgn}(\sigma)f(v_{\sigma(1)}, ..., v_{\sigma(n)}).\]
    \end{enumerate}
\end{definition}

\begin{definition}
    Niech \(V_1, ..., V_n = V\) to przestrzeń liniowa. Wówczas \(n\)-krotny iloczyn tensorowy
    \(V_1 \otimes ... \otimes V_n =: V^{\otimes_n}\) nazywamy \(n\)-tą potęgą
    tensorową przestrzeni \(V\).
\end{definition}

\begin{definition}
    Niech \(n \geq 1\) oraz \(V\) będzie przestrzenią liniową nad \(K\).
    Rozważmy następujące podprzestrzenie \(V^{\otimes_n}\):
    \begin{enumerate}
        \item \(N_s = \text{lin}\{v_1 \otimes ... \otimes v_n - v_{\sigma(1)}
        \otimes ... \otimes v_{\sigma(n)}\ :\ v_1, ..., v_n \in V \ \land \sigma \in S_n \}\);
        \item \(N_a = \text{lin}\{v_1 \otimes ... \otimes v_n\ :\ v_1, ..., v_n \in V \land v_i = v_j, i \not = j\}\)
    \end{enumerate}
    Zdefiniujmy \(\sum^{n}V = V^{\otimes} / N_s\) oraz \(\bigwedge^{n}V = V^{\otimes_n} / N_a\).
    Również niech
    \[
        \tau_s = p_s \circ \tau,\ \tau_a = p_a \circ \tau
    \]
    gdzie \(\tau: V^n \to V^{\otimes_n}\) to kanoniczne przekształcenie \(n\)-liniowe,
    a \(p_s: V^{\otimes_n} \to \sum^{n}V\) oraz \(p_a: V^{\otimes_n} \to \bigwedge^{n}V\)
    to naturalne rzutowania na przestrzenie ilorazowe. Wówczas:
    \begin{enumerate}
        \item przestrzeń \(\sum^{n}V\) nazywamy \(n\)-tą potęgą symetryczną \(V\);
        \item przestrzeń \(\bigwedge^{n}V\) nazywamy \(n\)-tą potęgą zewnętrzną.
        Elementy tej przestrzeni nazywamy czasem \(n\)-wektorami lub multiwektorami, a
        elementy \(\wedge^{n}V^{*}\) nazywamy \(n\)-formami lub \(n\)-kowektorami.
    \end{enumerate}
    Obraz \((v_1, ..., v_n)\) w \(p_a\) oznaczamy przez \(v_1 \wedge ... \wedge v_n\)
    i nazywamy \(n\)-wektorem (multiwektorem) prostym.
\end{definition}


\section{Iloczyn skalarny i ortogonalność}
\begin{definition}
    Formę dwuliniową \(\langle, \rangle\) nazywamy iloczynem skalarnym jeżeli:
    \begin{enumerate}
    \item Jest symetryczna;
    \item Jest dodatnio określona.
    \end{enumerate}
\end{definition}

\begin{definition}
    Definiujemy długość (normę) wektora \(\alpha\) jako liczbę
    \[ \norm{\alpha} = \sqrt{ \ip{\alpha}{\alpha} } \]
\end{definition}

\begin{lemma}[Podstawowe własności normy]
    Zachodzą następujące własności:
    \begin{enumerate}
        \item \( \norm{\alpha + \beta}^2 = \norm{\alpha}^2 + \norm{\beta}^2 + 2 \ip{\alpha}{\beta} \);
        \item \( \norm{a \cdot \alpha } = a \norm{\alpha}\);
        \item \( \norm{ \alpha } = 0 \Longleftrightarrow \alpha = 0 \).
    \end{enumerate}
\end{lemma}
\begin{lemma}[Uogólnienia nierówności szkolnych]
    Dla iloczynu skalarnego \(\ip{}{}\) prawdziwe są nierówności:
    \begin{enumerate}
        \item Nierówność Schwarza --- \(\abs{\ip{\alpha}{\beta}} \leq
        \norm{\alpha} \cdot \norm{\beta}\);
        \item Nierówność trójkąta --- \(\norm{\alpha + \beta} \leq \norm{\alpha}
        + \norm{\beta}\);
        \item Twierdzenie Pitagorasa --- \(\ip{\alpha}{\beta} = 0 \Longrightarrow
        \norm{\alpha + \beta}^2 = \norm{\alpha}^2 + \norm{\beta}^2\).
    \end{enumerate}
\end{lemma}

\begin{definition}
    Kątem między wektorami \(\alpha, \beta\) nazwiemy liczbę \(\theta\) taką, że:
    \[\text{cos}\theta = \frac{\ip{\alpha}{\beta}}{\norm{\alpha}\norm{\beta}}.\]
\end{definition}
\newpage
\begin{definition}
	Dla formy dwuliniowej \(f\) na \(V\) mówimy, że wektory \(\alpha, \beta\) są
	prostopadłe, gdy zachodzi \(f(\alpha, \beta) = 0\) i oznaczamy wtedy \(\alpha \perp \beta\).
	Definiujemy również dopełnienie ortogonalne podprzestrzeni \(W \subseteq V\)
	jako \[W^{\perp} = \{v \in V : \forall_{w \in W}\ v \perp w\}\]
\end{definition}

\begin{definition}
    Układ \((\alpha_1, ..., \alpha_n)\) wektorów nazwiemy ortogonalnym, jeśli każdej
    pary różnych \(i, j\) jest \(\alpha_i \perp \alpha_j\).
    Bazą ortagonalną nazwiemy układ który jest bazą i jest ortogonalny.
\end{definition}

\begin{remark}
    Zachodzą własności:
    \begin{enumerate}
        \item Jeśli \(X \subseteq V\) to ortogonalny układ niezerowych wektorów
        przestrzeni \(V\) z iloczynem skalarnym. Wówczas układ \(X\) jest
        liniowo niezależny.

        \item Każda przestrzeń euklidesowa ma bazę prostopadłą.
    \end{enumerate}
\end{remark}

\begin{definition}
	Mówimy, że układ jest ortonormalny, jeżeli jest ortogonalny
	oraz wszystkie jego wektory mają normę 1.
\end{definition}

\begin{remark}
    Prawdziwy jest następujący rozkład, gdzie pierwszy wyraz należy do
    \(\text{lin}(\alpha)\), a drugi do \(\text{lin}(\alpha)^{\perp}\):
    \[
    \beta = \frac{\ip{\alpha}{\beta}}{\ip{\alpha}{\alpha}}\alpha + \left( \beta
    - \frac{\ip{\alpha}{\beta}}{\ip{\alpha}{\alpha} \alpha} \right)
    \]
\end{remark}

\begin{remark}
    Niech \(V\) to \(n\)-wymiarowa przestrzeń liniowa z iloczynem skalarnym. Dla podprzestrzeni \(W \subseteq V\) zachodzi:
    \begin{enumerate}
        \item \(V = W \oplus W^{\perp}\);
        \item \(\text{dim}(W) + \text{dim}(W^{\perp}) = n\).
        \item \((W^{\perp})^{\perp} = W\)
    \end{enumerate}
\end{remark}

\begin{remark}
    Niech \(\mathcal{A} = (\alpha_1, ..., \alpha_n)\) to baza ortogonalna \(V\).
    Wówczas dla każdego wektora \(a_1 \alpha_1 + ... + a_n \alpha_n = \alpha \in
    V\) zachodzi
    \[
    a_1 = \frac{\ip{\alpha}{\alpha_1}}{\ip{\alpha_1}{\alpha_1}}, ..., a_n =
    \frac{\ip{\alpha}{\alpha_n}}{\ip{\alpha_n}{\alpha_n}}
    \]
\end{remark}


\begin{definition}
	Niech \(V\) to przestrzeń euklidesowa z podprzestrzenią \(W\). Wówczas:
	\begin{enumerate}
	    \item przekształczenie liniowe \(\phi: V \to V\) będące rzutem na \(W\)
			wzdłuż \(W^{\perp}\) nazywamy rzutem prostopadłym na \(W\).
		\item przekształcenie liniowe \(\psi: V \to V\) będące symetrią względem
		\(W\) wzdłuż \(W^{\perp}\) nazywamy symetrią prostopadłą względem \(W\).
    \end{enumerate}
\end{definition}

\begin{theorem}[Ortogonalizacja Grama-Smidta]
    Jesli \(V\) jest przestrzenią euklidesową oraz \((\gamma_1, ..., \gamma_n)\)
    jest bazą przestrzeni \(V\), to układ \((\alpha_1, ..., \alpha_n)\) zadany
    rekurencyjnie przez \(\alpha_1 = \gamma_1\) oraz dla \(k > 1\)
    \[
        \alpha_k = \gamma_k - \sum_{i = 1}^{k - 1}
        \frac{\ip{\gamma_k}{\alpha_i}}{\ip{\alpha_i}{\alpha_i}}
    \]
    jest bazą ortogonalną \(V\). Ponadto dla każdego \(k\) mamy
    \(\textnormal{lin}(\alpha_1, ..., \alpha_k) = \textnormal{lin}(\gamma_1, ..., \gamma_k)\).
\end{theorem}

\begin{remark}
    Z powyższego twierdzenia wynika kilka rzeczy:
    \begin{enumerate}
        \item Każdy układ ortogonalny mozna dopełnić do bazy ortogonalnej.
        \item Jeśli \(\phi \in \text{End}(V)\) jest triangularyzowalny, to
        istnieje też taka baza ortogonalna \(\mathcal{A}\) w której macierz
        \(\text{M}(\phi)_{\mathcal{A}}^{\mathcal{A}}\) jest górnotrójkątna.
    \end{enumerate}
\end{remark}

\begin{definition}
    Dla \(\)Wektor \(\alpha\) nazwiemy
\end{definition}

\section{Przekształcenia samosprzężone i twierdzenie spektralne}

Niech \(V, \ip{}{}\) to przestrzeń euklidesowa wymiaru \(n\). Każdemu wektorowi
\(v \in V\) przyporządkowujemy funkcjonał \(f_v \in V^{*}\) wzorem:
\[
    f_{v}(u) = \ip{u}{v},\ u \in V.
\]
Przyporządkowanie \(\Psi: V \to V^{*}\) zadane wzorem \(\Psi(v) = f_v\) jest
izomorfizmem przestrzeni liniowych. Jeśli \(\mathcal{A} = (\alpha_1, ..., \alpha_n)\)
jest bazą ortonormalną \(V\) to \((f_{v_1}, ..., f_{v_n})\) jest dualną do
\(\mathcal{A}\) bazą \(V^{*}\).
TODO

\section{Iloczyn hermitowski}

\begin{definition}
    Niech \(V\) to przestrzeń liniowa nad \(\C\). Funkcję \(\hip{}{} : V \times V \to \C\)
	nazywamy iloczynem hermitowskim, gdy dla każdych \(\alpha, \beta, \gamma, \delta \in V\)
	i \(a, b, c, d\) zachodzą:
	\begin{enumerate}
		\item \(\hip{a\alpha + b \beta}{\gamma} = a\hip{\alpha}{\gamma} + b\hip{\beta}{\gamma}\)
		\item \(\hip{\gamma}{a\alpha + b \beta} = \overline{a}\hip{\gamma}{\alpha} + \overline{b}\hip{\gamma}{\beta}\)
		\item \(\hip{\alpha}{\beta} = \overline{\hip{\beta}{\alpha}}\)
		\item \(\alpha \not = 0 \Rightarrow \hip{\alpha}{\alpha} \in \R_{+}\)
    \end{enumerate}
    Jeśli funkcja spełnia dwa pierwsze warunki, nazywamy ją formą półtoraliniową.
    Parę \((V, \hip{}{} )\), gdzie \(V\) jest skończenie wymiarowną przestrzenią liniową nac \(\C\)
    oraz \(\hip{}{} \) jest iloczynem hermitowskim na \(V\) nazywamy przestrzenią unitarną.
\end{definition}

\begin{remark}
	Dla iloczynu hermitowskiego działają również macierze form dwuliniowych.
\end{remark}

\begin{remark}
	Definicje ortogonalności, normy i baz tak jak w iloczynie skalarnym.
\end{remark}
\begin{remark}
	Jeśli \(\alpha, \beta \in V\) oraz \(\hip{\alpha}{\beta} = 0\) wtw. gdy
	\(\hip{\beta}{\alpha} = 0\).
\end{remark}

\begin{remark}
	Nierówności Schwarza, trójkąta i twierdzenie Pitagorasa działają również
	w iloczynie hermitowskim.
\end{remark}

\begin{definition}
	Niech \(A = [a_{i\,j}]\). Macierz \(A^{*}\) mającą w \(i\)-tym wierszu i \(j\)-tej
	kolumnie wyraz \(\overline{a_{j\,i}}\) nazywamy sprzężeniem hermitowskim
	macierzy \(A\). Jeśli macierz \(A\) spełnia warunek \(A^{*} = A\), to
	mówimy, że \(A\) jest macierzą hermitowską.
\end{definition}

\begin{definition}
	Endomorfizm \(\phi\) nazywamy samosprzężonym jeśli \(\phi = \phi^{*}\).
\end{definition}

\begin{definition}
	Niech \((V_1, \hip{}{}_1), (V_2, \hip{}{}_2)\) to przestrzenie unitarne.
	Wówczas:
	\begin{enumerate}
		\item Izomorfizm \(\phi: V_1 \to V_2\) nazwiemy przekształceniem unitarnym, jeżeli
		zachowuje iloczyn hermitowski, czyli \(\hip{\alpha}{\beta}_1 = \hip{\phi(\alpha)}{\phi(beta)}_2\);

		\item Macierz \(A\) nazwiemy unitarną jeśli \(M \cdot M^{*} = I_n\).
    \end{enumerate}
\end{definition}

\begin{remark}
	Izomorfizm jest unitarny wtw. gdy zachowuje normę każdego wektora \(v \in V_1\).
\end{remark}

\begin{theorem}
	Jeśli \(\phi \in \textnormal{End}(V)\) jest endomorfizmem unitarnym, to
	każda wartość własna \(\lambda \in \C\) spełnia \(\abs{\lambda} = 1\).
	Wektory własne \(\phi\) odpowiadające parami różnym wartościom własnym są
	ortogonalne. Istnieje też baza złożona z wektorów własnych.
\end{theorem}

\section{SVD}


\end{document}
